\documentclass[11pt]{article}

\usepackage[letterpaper,margin=1in]{geometry}
\usepackage{amsmath,amssymb,mathtools}
\usepackage{booktabs}
\usepackage{tabularx}
\usepackage{enumitem}
\usepackage{parskip}
\usepackage{microtype}
\usepackage[hidelinks]{hyperref}

% ---- numbered, continuable definition list ----
\newlist{definitionlist}{enumerate}{1}
\setlist[definitionlist]{label=\textbf{(\arabic*)}, leftmargin=2.6em,
  labelsep=0.6em, itemsep=0.5ex, topsep=0.7ex}

\title{A Formal Definition of Instructions in a\\[2pt]
       Domain-Specific Language}
\author{Prepared with Claude (Anthropic)}
\date{\today}

\begin{document}
\maketitle

\begin{abstract}
An \emph{instruction} (or \emph{command}) sits at the intersection of two things:
a symbol in the grammar of a domain-specific language, and a behavior in the
program that language describes. This note defines it self-containedly -- its
abstract syntax as a named operator carrying optional parameters, the forms those
parameters take, its reading across the standard grammar metasyntaxes, and its
semantics as a function from arguments to a computation -- with no dependence on
any external apparatus.
\end{abstract}

\section*{Instructions}

An \textbf{instruction} (equivalently, a \textbf{command}) is a node of the
abstract syntax: a named operator that may carry data. Its name is the node's
head; its data are the node's children. No meaning resides in the name itself --
meaning is supplied separately, by an interpretation, developed in the final
section.

Given:

\begin{definitionlist}
    \item a \textbf{name} \(n\), an identifier drawn from a namespace
          \(\mathcal{N}\);
    \item a \textbf{parameter} \(a\): a slot the instruction may carry, equipped
          with a \textbf{domain} \(D_a\) (its admissible values) and a marking
          \emph{required} or \emph{optional}; an optional parameter additionally
          admits the \textbf{absent} value \(\bot \notin D_a\) and may name a
          \textbf{default} \(d_a \in D_a\);
    \item a \textbf{parameter profile} \(\operatorname{prof}(n)\): the finite,
          possibly empty, set of parameters declared for the name \(n\),
\end{definitionlist}

one defines:

\begin{definitionlist}[start=4]
    \item an \textbf{instruction} \(c\): a name \(\operatorname{name}(c)\in
          \mathcal{N}\) together with a \textbf{binding} \(\beta_c\) assigning to
          each \(a\in\operatorname{prof}(c)\) a value \(\beta_c(a)\in
          D_a\cup\{\bot\}\), with \(\bot\) permitted only where \(a\) is
          optional. The \textbf{arity} of \(c\) is
          \(\lvert\operatorname{prof}(c)\rvert\), writing
          \(\operatorname{prof}(c):=\operatorname{prof}(\operatorname{name}(c))\).
\end{definitionlist}

The binding is the instruction's entire variable content: a finite map from its
parameters to values, recording which optional parameters are present and what
each carries. Flags, attributes, arguments -- and command-line ``options'' in the
informal sense -- are not separate primitives but \emph{parameters} differing only
in the shape of their domain; an instruction ``may or may not'' carry them exactly
as its profile is non-empty or empty. An instruction of empty profile is a
\textbf{bare command}, an atom with a name and nothing more. Being \emph{abstract}
syntax, an instruction retains only its name and binding; distinct surface
renderings that induce the same binding are the same instruction.

Each parameter falls, by the shape of \(D_a\), into exactly one \textbf{mode}:

\begin{definitionlist}[start=5]
    \item \textbf{nullary} (\emph{flag}): \(D_a=\{\star\}\), a single presence
          token; the only datum is present-versus-absent.
    \item \textbf{Boolean}: \(D_a=\mathbb{B}=\{\mathsf{true},\mathsf{false}\}\).
    \item \textbf{enumerated}: \(D_a=\{u_1,\dots,u_r\}\), a finite enumeration.
    \item \textbf{valued}: \(D_a\) is an arbitrary type (the general case).
\end{definitionlist}

These nest -- nullary is the one-element enumeration, Boolean the two-element one
-- and are separated only so the semantics can treat their values uniformly.

\section*{The forms an instruction takes}

The variable content of an instruction is the values its binding records. A few
recurring forms are worth naming; each is a statement about the binding alone.

\begin{table}[h]
\centering
\footnotesize
\setlength{\tabcolsep}{6pt}
\renewcommand{\arraystretch}{1.35}
\begin{tabularx}{\textwidth}{@{}l l >{\raggedright\arraybackslash}X@{}}
\toprule
\textbf{Surface instruction} & \textbf{Domain \(D_a\)} &
\textbf{Binding contribution} \(\bigl(\text{parameter}\mapsto\text{value}\bigr)\)\\
\midrule
flag \(\mathsf{foo}\) (present) & \(\{\star\}\) &
\(\mathsf{foo}\mapsto\star\) \;-- present, valueless\\
\(\mathsf{do\_foo}=b\) & \(\mathbb{B}\) &
\(\mathsf{do\_foo}\mapsto b,\quad b\in\{\mathsf{true},\mathsf{false}\}\)\\
pair \(\mathsf{do\_foo}\,/\,\mathsf{dont\_do\_foo}\) & name-selected \(\mathbb{B}\) &
\(s\mapsto\varpi(\operatorname{name}(c))\) \;for a shared setting \(s\)\\
\(\mathsf{some\_command}=u\) & enum \(\{u_1,\dots,u_r\}\) &
\(\mathsf{some\_command}\mapsto u,\quad u\in\{u_1,\dots,u_r\}\)\\
\bottomrule
\end{tabularx}
\caption{Each instruction form as a contribution to the binding, graded by
\(D_a\).}
\end{table}

The \textbf{flag} \(\mathsf{foo}\) carries a single nullary parameter: issued, its
binding records presence with no value; unissued, the parameter is absent
(\(\bot\)). Only presence is observable, and presence \emph{is} the content.

The \textbf{Boolean} \(\mathsf{do\_foo}=b\) records a truth value directly.

The \textbf{polarity pair} is two distinct command names, neither carrying an
explicit value, that denote one underlying Boolean setting \(s\) with opposite
values. A \textbf{polarity} \(\varpi\) sends each name to a Boolean --
\(\varpi(\mathsf{do\_foo})=\mathsf{true}\),
\(\varpi(\mathsf{dont\_do\_foo})=\mathsf{false}\) -- so the choice of name selects
the value of \(s\). The pair \(\mathsf{do\_foo}/\mathsf{dont\_do\_foo}\) and the
explicit form \(\mathsf{do\_foo}=b\) therefore carry the \emph{same} content: two
concrete syntaxes for one setting. (The two names form an orbit of the
fixed-point-free involution that swaps polarity, and \(\varpi\) is the
\(\mathbb{Z}/2\)-valued invariant it induces.)

The \textbf{enumerated} \(\mathsf{some\_command}=u\) records a value from a finite
domain.

\textbf{Remark (one scheme).} The forms are a single construction graded by
\(\lvert D_a\rvert\) and by whether the value is retained or suppressed to
presence alone: the flag is the one-element domain with its value forgotten, the
Boolean the two-element domain, the enumeration the finite domain, the valued
parameter the unrestricted one. A multi-parameter instruction is just their
product over the profile.

\section*{Reading the definition against a grammar}

The construction does not depend on how a DSL's surface is written down. Fix a
context-free grammar \(G=(N,\Sigma,P,S)\) -- nonterminals \(N\), terminals
\(\Sigma\), productions \(P\), start symbol \(S\) -- presented in any of the usual
notations. The instruction layer is the \emph{abstract} content of \(G\); the
notation is only its concrete surface. The dictionary is:

\begin{definitionlist}
    \item a \textbf{name} is a terminal keyword in \(\Sigma\) (equivalently, the
          production the keyword selects);
    \item an \textbf{instruction} is a nonterminal \(\langle\textit{cmd}\rangle\)
          together with its production; the \textbf{profile} is that production's
          right-hand side;
    \item a \textbf{parameter} \(a\) is a nonterminal \(A_a\), and its
          \textbf{domain} is the language that nonterminal generates,
          \(D_a=\mathcal{L}(A_a)\);
    \item a \textbf{binding} \(\beta_c\) is the attribution of a parse tree: a
          filled parameter slot carries the value parsed there; an optional slot
          taken empty carries \(\bot\);
    \item the \textbf{mode} is read off \(A_a\)'s production -- a singleton
          right-hand side (nullary), \(\mathsf{true}\mid\mathsf{false}\)
          (Boolean), a finite alternation (enumerated), or an arbitrary
          sublanguage (valued).
\end{definitionlist}

Taking \(D_a=\mathcal{L}(A_a)\) is what frees the definition from any particular
DSL: each language fixes its own parameter languages, and the rest of the
development is identical over all of them.

\textbf{The notations differ only in surface.} Each metasyntax expresses the same
operators, so one production translates uniformly:

\begin{table}[h]
\centering
\footnotesize
\setlength{\tabcolsep}{6pt}
\renewcommand{\arraystretch}{1.3}
\begin{tabularx}{\textwidth}{@{}l *{4}{>{\raggedright\arraybackslash}X}@{}}
\toprule
 & \textbf{BNF} & \textbf{EBNF} & \textbf{ABNF} & \textbf{PEG}\\
\midrule
definition          & \texttt{::=} & \texttt{=} & \texttt{=} & \(\leftarrow\)\\
alternation         & \texttt{|} & \texttt{|} & \texttt{/} & \texttt{/} \emph{(ordered)}\\
optional \((X)\)    & \(X\mid\varepsilon\) & \texttt{[X]} & \texttt{[X]} & \texttt{X?}\\
repetition \((0{+})\) & recursion & \texttt{\{X\}} & \texttt{*X} & \texttt{X*}\\
terminal            & literal & \texttt{"lit"} & \texttt{\%xHH}, \texttt{"lit"} & \texttt{"lit"}, \texttt{[class]}\\
nonterminal         & \(\langle\textit{name}\rangle\) & \textit{name} & \textit{rulename} & \textit{Name}\\
\bottomrule
\end{tabularx}
\caption{One production across the four metasyntaxes. Alternation is
\emph{unordered} in BNF, EBNF, and ABNF -- a symmetric choice whose ambiguity is
resolved outside the grammar -- but \emph{ordered} in PEG, where \texttt{/} is a
committed, first-match choice.}
\end{table}

\textbf{Presence as a single bit.} An instruction ``whose presence or absence is
enough'' is exactly the optionality operator: \(X\mid\varepsilon\) in BNF,
\([X]\) in EBNF and ABNF, \(X\texttt{?}\) in PEG. Its parameter is nullary, or a
parameter whose value is discarded, so its binding records only presence; the
syntactic domain may be a singleton or richer, but the observable content is the
single bit present-versus-absent.

\textbf{Precedence.} Priority enters at two independent layers.

\emph{Syntactic.} When several alternatives could apply, a choice must be made.
PEG fixes it inside the grammar: \(e_1/e_2\) tries the left alternative first and
commits on success, so precedence is the order of alternatives and the grammar is
unambiguous by construction. BNF, EBNF, and ABNF leave alternation symmetric
(ABNF's \texttt{/} is set union, not PEG's ordered choice, despite the shared
glyph), so precedence and associativity are imposed from outside -- by
stratifying the grammar into levels
\(\bigl(\langle\textit{expr}\rangle,\langle\textit{term}\rangle,
\langle\textit{factor}\rangle\bigr)\) or by separate precedence declarations.

\emph{Semantic.} When a program issues conflicting settings for one parameter, a
resolution policy decides the winner; the usual policy is sequential override -- a
later instruction supersedes an earlier one. The two layers are dual instances of
one device, a total order that makes a nondeterministic combine deterministic:
PEG's ordered choice is ``leftmost alternative wins'' at parse time, sequential
override is ``latest issuance wins'' at execution time.

\textbf{Attributes.} The denotation is an \textbf{attribute grammar} over \(G\).
Each parameter nonterminal \(A_a\) \emph{synthesizes} its value
\(\beta_c(a)\in D_a\cup\{\bot\}\); the production for \(\langle\textit{cmd}\rangle\)
carries a semantic rule assembling these into the instruction's binding, and a
further rule into its behavior. Defaults \(d_a\) and any ambient context -- the
active namespace, the enclosing scope -- travel the other way, as \emph{inherited}
attributes. Because every value is synthesized and the only inheritance flows
left-to-right, the definition is \(L\)-attributed: it is evaluable in a single
left-to-right pass.

\section*{The behavioral half: instructions as interpreted symbols}

The previous sections fixed the \emph{syntactic} half of an instruction -- a
name, a profile, and the binding it parses to. Its \emph{semantic} half is a
function \(f\) saying what executing it does. The two are glued along the profile:
parsing fills the profile with a binding, and \(f\) reads that binding. An
instruction is thus one symbol wearing two faces -- a grammar symbol and a
behavior -- paired over the arity they share.

Fix:

\begin{definitionlist}
    \item an \textbf{effect monad} \(T\) and a \textbf{result type} \(R\); a
          \textbf{computation} is an element of \(T(R)\). The common command runs
          only for effect, \(R=\mathbf{1}\); specializing \(T\) to the State monad
          \(T(R)=\mathit{St}\to(R\times\mathit{St})\) makes a computation a state
          transformer \(\mathit{St}\to\mathit{St}\), while other monads supply
          exceptions, input/output, or nondeterminism;
    \item the \textbf{binding space}, the dependent sum (a telescope in which each
          factor may depend on the earlier ones)
          \[
            B \;=\; \sum_{a\in A}\bigl(D_a(\beta_{<a})\cup\{\bot\}\bigr),
          \]
          degenerating to the flat product \(\prod_{a\in A}(D_a\cup\{\bot\})\) when
          the parameter domains are independent;
    \item a \textbf{validity predicate} \(\Phi\colon B\to\mathbb{B}\), the
          conjunction of the cross-parameter constraints -- required-present,
          mutual exclusion, \(\textsf{requires}\)/\(\textsf{implies}\) -- cutting
          out the \textbf{valid bindings}
          \[
            B_{\mathrm{valid}} \;=\; \{\,\beta\in B \mid \Phi(\beta)\,\}.
          \]
\end{definitionlist}

Dependencies -- one parameter's \emph{domain} gated by another's value -- live in
the \(\Sigma\)-shape of \(B\); flat constraints -- exclusion, co-requirement --
live in \(\Phi\).

\textbf{The semantic function.} The behavior of an instruction is a map

\[
  f\colon B_{\mathrm{valid}} \longrightarrow T(R),
\]

total on the valid bindings. Presence or absence of a flag is the choice of
\(\bot\) in a coordinate, so \(f\) is a single function on a (dependent) product
-- not one function per subset of present flags. The space of candidate behaviors
compatible with the symbol's arity is therefore the exponential

\[
  F \;=\; \bigl[\,B_{\mathrm{valid}}\to T(R)\,\bigr],
\]

and an instruction is its symbol together with a chosen point \(f\in F\).

\textbf{Assembling the behavior parameter-wise.} Defining \(f\) pointwise would
list \(\lvert B_{\mathrm{valid}}\rvert\) cases; present it instead from parts
linear in the parameters. Give the instruction a \textbf{base behavior} \(f_0\), a
function of its required arguments alone, and give each optional parameter \(a\) a
\textbf{modifier} \(m_a\colon D_a\cup\{\bot\}\to\operatorname{End}\bigl(T(R)\bigr)\),
an endo-transformer on computations, with

\[
  m_a(\bot)=\mathrm{id}.
\]

Then

\[
  f(\beta) \;=\;
  \Bigl(\,\mathop{\bigcirc}\limits_{a\in A_{\mathrm{opt}}} m_a(\beta(a))\,\Bigr)
  \bigl(f_0(\beta|_{\mathrm{req}})\bigr),
\]

the composite of the present modifiers applied to the base, where
\(A_{\mathrm{opt}}\) is the set of optional parameters and
\(\beta|_{\mathrm{req}}\) the restriction of \(\beta\) to the required ones.
Absence is the identity element, presence applies a modifier, and a combination is
their composite: \(\lvert A_{\mathrm{opt}}\rvert\) modifiers and one base generate
every behavior, the exponential blow-up confined to the \emph{image}
\(\{\,f(\beta):\beta\in B_{\mathrm{valid}}\,\}\) -- the orbit of \(f_0\) under the
modifier monoid -- which is never written out.

\textbf{Proposition (order-independence \(\iff\) commutativity).} The collection
of present optional parameters carries no intrinsic order, so the composite above
is well defined exactly when the modifiers pairwise commute in
\(\operatorname{End}(T(R))\); equivalently, the behavior factors through the action
of a \emph{commutative} monoid, and parameter order is then semantically
irrelevant. When the modifiers do not commute, the composite is made well defined
by composing them in the parameter precedence of the preceding section.

\emph{Proof sketch.} A fold over an unordered set is independent of its
enumeration iff its operations commute, which for the monoid
\((\operatorname{End}(T(R)),\circ,\mathrm{id})\) is pairwise commutativity of the
\(m_a(\beta(a))\). Failing that, any total order on the parameters determines the
fold, and precedence is the canonical such order. \(\square\)

\textbf{Sequencing.} A \textbf{program} is a finite sequence \(c_1;\dots;c_m\) of
instructions. Its behavior is the composite of the individual computations in the
effect monad,

\[
  f_{c_m}(\beta_{c_m}) \;\bullet\; \cdots \;\bullet\; f_{c_1}(\beta_{c_1}),
\]

where \(\bullet\) is composition in \(T\) -- for the State monad, ordinary
composition of state transformers, threading the state through -- applied right to
left, so \(c_1\) executes first. Sequential override is then a derived property:
if two instructions write the same state component, the later computation runs
second and its write prevails.

\textbf{The intersection, precisely.} Write \(\Omega\) for the signature of
instruction constructors. The abstract syntax of the first section is the term
algebra it freely generates; assigning every constructor a behavior makes
computations an \(\Omega\)-algebra, and the program's denotation is the unique
homomorphism from that term algebra into the algebra of computations. The function
\(f\) is this homomorphism's action at a single constructor. An instruction is
therefore one constructor read in both algebras at once -- a symbol on the
syntactic side, a computation on the semantic side, identified along the profile
each must honor. That single arrow, from a symbol to its computation, is the
intersection of grammar and behavior the development has been circling.

\end{document}
